%! AP Statistics — Unit 2, Lesson 2.8 — Group Activity — Answer Key
\documentclass[10pt]{article}
\usepackage{apstats-article}
\usepackage{apstats-key}

\begin{document}

\pageheader{Unit 2, Lesson 2.8}{Group Activity --- Answer Key}
\namepartnerperiod

\vspace{0.12in}

\begin{tcolorbox}[colback=sky, colframe=navy, boxrule=0.9pt, arc=2mm,
    left=4mm, right=4mm, top=2.5mm, bottom=2.5mm]
\small\textit{A school nutritionist collects data on 40 students. For each student she records
daily screen time in hours ($x$) and a composite wellness score ($y$, 0--100 scale, higher is
better). Summary statistics: $\bar{x} = 4.2$ hr, $\bar{y} = 71$, $s_x = 1.8$ hr, $s_y = 12$,
$r = -0.75$.}
\end{tcolorbox}

\vspace{0.14in}

% ── TIER R ───────────────────────────────────────────────────────────────────
\begin{tcolorbox}[colback=redbg, colframe=black!40, boxrule=0.7pt, arc=2mm,
    left=4mm, right=4mm, top=2.5mm, bottom=2.5mm,
    title={\textbf{Tier R --- Remediate}}, fonttitle=\bfseries]
\small

\textbf{Use the formula reference below for each step.}

\vspace{4pt}
\begin{tabularx}{\linewidth}{@{} >{\bfseries}p{4cm} X @{}}
Slope formula & $b = r \cdot \dfrac{s_y}{s_x}$ \\[4pt]
Intercept formula & $a = \bar{y} - b\bar{x}$ \\
\end{tabularx}

\vspace{8pt}
\begin{enumerate}[leftmargin=1.6em, itemsep=0.16in]
  \item Substitute the values into the slope formula and compute $b$:\\[6pt]
        \ansline{$b = (-0.75) \cdot \dfrac{12}{1.8} = -5$ points per hour}

  \item Substitute $b$, $\bar{x}$, and $\bar{y}$ into the intercept formula and compute $a$:\\[6pt]
        \ansline{$a = 71 - (-5)(4.2) = 71 + 21 = 92$}

  \item Write the regression equation using the variable names (not $x$ and $y$):\\[6pt]
        \ansline{$\widehat{\text{wellness}} = 92 - 5(\text{screen time})$}

  \item Compute $r^2$: \ans{$(-0.75)^2 = 0.5625$}. Circle the correct interpretation:

        \vspace{6pt}
        \textbf{(A)} About 56\% of students have a wellness score that is linear.\\
        \textbf{(B)} About 56\% of the variation in wellness score is explained by the regression on screen time.
        \textcolor{keyred}{\textbf{$\leftarrow$ correct}}\\
        \textbf{(C)} The correlation between screen time and wellness is 56\%.
\end{enumerate}
\end{tcolorbox}

\vspace{0.12in}

% ── TIER A ───────────────────────────────────────────────────────────────────
\begin{tcolorbox}[colback=goldbg, colframe=black!40, boxrule=0.7pt, arc=2mm,
    left=4mm, right=4mm, top=2.5mm, bottom=2.5mm,
    title={\textbf{Tier A --- Approaching Proficiency}}, fonttitle=\bfseries]
\small

\begin{enumerate}[leftmargin=1.6em, itemsep=0.16in]
  \item Compute $b$ and $a$ from the summary statistics. Show all work.\\[8pt]
        \ansline{$b = (-0.75)\cdot\frac{12}{1.8} = -5$ points per hour}\\[3pt]
        \ansline{$a = 71 - (-5)(4.2) = 71 + 21 = 92$}

  \item Write the regression equation with variable names. Then interpret the slope in context.\\[8pt]
        \ansline{$\widehat{\text{wellness}} = 92 - 5(\text{screen time})$}\\[3pt]
        \ansline{For each additional hour of daily screen time, the predicted wellness score}\\[3pt]
        \ansline{decreases by approximately 5 points.}

  \item Compute $r^2$ and write the full AP-template interpretation.\\[8pt]
        \ansline{$r^2 = (-0.75)^2 = 0.5625$}\\[3pt]
        \ansline{About 56.25\% of the variation in wellness score is explained by the}\\[3pt]
        \ansline{least-squares regression of wellness score on daily screen time.}

  \item Predict the wellness score for a student with 5 hours of daily screen time. Classify
        as interpolation or extrapolation. (The data ranged from 1 to 9 hours.)\\[8pt]
        \ansline{$\widehat{\text{wellness}} = 92 - 5(5) = 92 - 25 = 67$ points.}\\[3pt]
        \ansline{\textbf{Interpolation} --- 5 hours is within the data range (1--9 hr).}
\end{enumerate}
\end{tcolorbox}

\vspace{0.12in}

% ── TIER E ───────────────────────────────────────────────────────────────────
\begin{tcolorbox}[colback=greenbg, colframe=black!40, boxrule=0.7pt, arc=2mm,
    left=4mm, right=4mm, top=2.5mm, bottom=2.5mm,
    title={\textbf{Tier E --- Extension}}, fonttitle=\bfseries]
\small

\begin{enumerate}[leftmargin=1.6em, itemsep=0.16in]
  \item Complete all of Tier A above first.

  \item A new data point is added: a student with 14 hours of daily screen time and a wellness
        score of 85. Is this point likely to be \textit{influential}? Explain using the concepts
        of leverage and trend.\\[8pt]
        \ansline{Yes. $x = 14$ hr is far outside the data range (1--9 hr), giving it high leverage.}\\[3pt]
        \ansline{Additionally, a wellness score of 85 at 14 hours goes \emph{against} the negative}\\[3pt]
        \ansline{trend, so it would pull the slope upward --- making it influential.}

  \item Without recomputing, predict whether adding this point would \textit{increase} or
        \textit{decrease} the slope, the intercept, and $r^2$. Justify each.\\[8pt]
        Slope: \ansline{Increase (become less negative) --- the point at high $x$ with high $y$ pulls the line up.}\\[3pt]
        Intercept: \ansline{Decrease --- as slope rises (becomes less steep), the intercept shifts down to keep the line through $(\bar{x},\bar{y})$.}\\[3pt]
        $r^2$: \ansline{Decrease --- the new point does not follow the negative trend, weakening the linear}\\[3pt]
        \ansline{relationship and reducing the proportion of variation explained.}

  \item The nutritionist removes the point and recomputes: $r^2$ changes from 0.5625 to 0.48.
        Does this confirm or contradict your prediction? What does the change tell us about the
        point's influence?\\[8pt]
        \ansline{This \textbf{confirms} the prediction: $r^2$ decreased when the point was present (0.48 without}\\[3pt]
        \ansline{it vs.\ 0.5625 with it would contradict --- re-read the problem: $r^2$ drops from 0.5625 to 0.48}\\[3pt]
        \ansline{when the point is removed, meaning the outlier \emph{inflated} $r^2$. Investigate context.}
\end{enumerate}
\end{tcolorbox}

\begin{teachernote}
Tier E item 3 is intentionally subtle: the outlier at high $x$ with high $y$ can sometimes
\emph{increase} $r^2$ if it extends the range in a way consistent with the trend, even though
it contradicts the general direction. Use this to push students beyond the simple rule ``outliers
always hurt the model'' and toward careful case analysis. The key takeaway: always investigate,
never automatically remove.
\end{teachernote}

\end{document}
